% Ngân hàng tự sinh (AI)
\begin{bt}%[0H9N1-1]
% nguon: AI soạn tương tự câu mẫu 959cabbe
% mau: 959cabbe
	Trên trục tọa độ $\left(O;\overrightarrow{i}\right)$ cho 2 điểm $A$, $B$ có tọa độ lần lượt là $-5; 3$. Tìm tọa độ của véc-tơ $\overrightarrow{AB}$.
	\loigiai{
		Ta có $\overrightarrow{AB}=3+5=8$.\\
		Vậy toạ độ $\overrightarrow{AB}$ là $8$.
	}
\end{bt}

\begin{bt}%[0H9H1-1]
% nguon: AI soạn tương tự câu mẫu 959cabbe
% mau: 959cabbe nang_cao
Trên trục tọa độ $\left(O;\overrightarrow{i}\right)$ cho hai điểm $M$, $N$ có tọa độ lần lượt là $-7$ và $4$. Gọi $P$ là điểm có tọa độ $-1$. Tính tọa độ của véc-tơ $\overrightarrow{MN}$ và tọa độ của véc-tơ $\overrightarrow{PN}$.
	\loigiai{
		Ta có $\overrightarrow{MN}=4-(-7)=4+7=11$. Vậy tọa độ $\overrightarrow{MN}$ là $11$.\\
		Ta có $\overrightarrow{PN}=4-(-1)=4+1=5$. Vậy tọa độ $\overrightarrow{PN}$ là $5$.
	}
\end{bt}

\begin{bt}%[0H9N1-2]
% nguon: AI soạn tương tự câu mẫu ae5e3c24
% mau: ae5e3c24
	Trong mặt phẳng tọa độ $Oxy$, cho $2$ véc-tơ $\overrightarrow{a}=(5;3)$, $\overrightarrow{b}=(-6;10)$. Tính tích vô hướng $\overrightarrow{a}\cdot \overrightarrow{b}$ và chứng tỏ $\overrightarrow{a}\perp\overrightarrow{b}$.
	\loigiai{
		Ta có $\overrightarrow{a}\cdot \overrightarrow{b} = 5\cdot (-6) + 3\cdot 10 =0$.\\
		Suy ra $\overrightarrow{a}\perp\overrightarrow{b}$.
	}
\end{bt}

\begin{bt}%[0H9H1-2]
% nguon: AI soạn tương tự câu mẫu ae5e3c24
% mau: ae5e3c24 nang_cao
	Trong mặt phẳng tọa độ $Oxy$, cho hai véc-tơ $\overrightarrow{a}=(6;-2)$ và $\overrightarrow{b}=(3;9)$. Tính tích vô hướng $\overrightarrow{a}\cdot \overrightarrow{b}$ và chứng tỏ $\overrightarrow{a}\perp\overrightarrow{b}$.
	\loigiai{
		Ta có $\overrightarrow{a}\cdot \overrightarrow{b} = 6\cdot 3 + (-2)\cdot 9 = 18-18 = 0$.\\
		Suy ra $\overrightarrow{a}\perp\overrightarrow{b}$.
	}
\end{bt}

\begin{bt}%[0H9N1-3]
% nguon: AI soạn tương tự câu mẫu 04017960
% mau: 04017960
Trong mặt phẳng tọa độ $Oxy$, cho hai điểm $P(-5; 2)$ và $Q(4; 6)$. Tìm tọa độ $\overrightarrow{PQ}$.
	\loigiai{
		Ta có $\overrightarrow{PQ}=(9; 4)$.
	}
\end{bt}

\begin{bt}%[0H9H1-3]
% nguon: AI soạn tương tự câu mẫu 04017960
% mau: 04017960 nang_cao
	Trong mặt phẳng tọa độ $Oxy$, cho hai điểm $A(-6; 3)$ và $B(5; -4)$. Tìm tọa độ $\overrightarrow{AB}$ và tính độ dài đoạn thẳng $AB$.
	\loigiai{
		Ta có $\overrightarrow{AB}=(5-(-6);\, -4-3)=(11; -7)$.\\
		Độ dài đoạn thẳng $AB$ là $AB=\sqrt{11^2+(-7)^2}=\sqrt{121+49}=\sqrt{170}$.
	}
\end{bt}

\begin{bt}%[0H9H1-4]
% nguon: AI soạn tương tự câu mẫu 2f967af4
% mau: 2f967af4
Trong mặt phẳng $Oxy$, cho $A\left(2; -1\right), B\left(4; 1\right), C\left(1; 3\right)$. Gọi $M$ là điểm thuộc $BC$ sao cho $\overrightarrow{MB}=-2 \overrightarrow{BC}$ và $G$ là trọng tâm của tam giác $ABC$.
    \begin{enumerate}
        \item Phân tích $\overrightarrow{AM}$ theo hai véc-tơ $\overrightarrow{AB}$ và $\overrightarrow{AC}$. 
        \item Phân tích $\overrightarrow{GM}$ theo hai véc-tơ $\overrightarrow{AB}$ và $\overrightarrow{AC}$.
    \end{enumerate}
	\loigiai{
\begin{enumerate}
        \item Gọi $M(x;y) \Rightarrow \overrightarrow{MB}=(4-x;1-y); \overrightarrow{BC}=(-3;2); \overrightarrow{AB}=(2;2); \overrightarrow{AC}=(-1;4)$.\\
        Mà $\overrightarrow{MB}=-2 \overrightarrow{BC} \Leftrightarrow \heva{&4-x=6\\ &1-y=-4} \Leftrightarrow \heva{&x=-2\\ &y=5} \Rightarrow M(-2;5) \Rightarrow \overrightarrow{AM}=(-4;6)$.\\
        Giả sử $\overrightarrow{AM} = m\overrightarrow{AB}+n\overrightarrow{AC} \Leftrightarrow \heva{&-4=2m-n\\ &6=2m+4n} \Leftrightarrow \heva{&m= -1\\&n=2.}$\\
        Vậy $\overrightarrow{AM} = -\overrightarrow{AB}+2\overrightarrow{AC}$.
        \item Ta có tọa độ trọng tâm $G\left(\dfrac{7}{3};1\right) \Rightarrow \overrightarrow{GM}=\left(-\dfrac{13}{3};4\right)$.\\
        Giả sử $\overrightarrow{GM} = h\overrightarrow{AB}+k\overrightarrow{AC} \Leftrightarrow \heva{&-\dfrac{13}{3}=2h-k\\ &4=2h+4k} \Leftrightarrow \heva{&h= -\dfrac{4}{3}\\ &k= \dfrac{5}{3}.}$\\
        Vậy $\overrightarrow{GM} = -\dfrac{4}{3}\overrightarrow{AB}+\dfrac{5}{3}\overrightarrow{AC}$.
    \end{enumerate} 
    \textit{\textbf{Chú ý:} Chúng ta hoàn toàn có thể sử dụng cách biến đổi véc-tơ để thực hiện việc biểu diễn $\overrightarrow{AM},\overrightarrow{GM}$ theo hai véc-tơ $\overrightarrow{AB}$ và $\overrightarrow{AC}$.}
	}
\end{bt}

\begin{bt}%[0H9V1-4]
% nguon: AI soạn tương tự câu mẫu 2f967af4
% mau: 2f967af4 nang_cao
Trong mặt phẳng $Oxy$, cho $A\left(1; -2\right), B\left(3; 0\right), C\left(-1; 4\right)$. Gọi $M$ là điểm thuộc $BC$ sao cho $\overrightarrow{MB}=-3\overrightarrow{BC}$ và $G$ là trọng tâm của tam giác $ABC$.
    \begin{enumerate}
        \item Phân tích $\overrightarrow{AM}$ theo hai véc-tơ $\overrightarrow{AB}$ và $\overrightarrow{AC}$.
        \item Phân tích $\overrightarrow{GM}$ theo hai véc-tơ $\overrightarrow{AB}$ và $\overrightarrow{AC}$.
    \end{enumerate}
	\loigiai{
		\begin{enumerate}
        \item Gọi $M(x;y) \Rightarrow \overrightarrow{MB}=(3-x;-y); \overrightarrow{BC}=(-4;4); \overrightarrow{AB}=(2;2); \overrightarrow{AC}=(-2;6)$.\\
        Mà $\overrightarrow{MB}=-3\overrightarrow{BC} \Leftrightarrow \heva{&3-x=12\\ &-y=-12} \Leftrightarrow \heva{&x=-9\\ &y=12} \Rightarrow M(-9;12) \Rightarrow \overrightarrow{AM}=(-10;14)$.\\
        Giả sử $\overrightarrow{AM} = m\overrightarrow{AB}+n\overrightarrow{AC} \Leftrightarrow \heva{&-10=2m-2n\\ &14=2m+6n} \Leftrightarrow \heva{&m= -2\\&n=3.}$\\
        Vậy $\overrightarrow{AM} = -2\overrightarrow{AB}+3\overrightarrow{AC}$.
        \item Ta có tọa độ trọng tâm $G(1;1) \Rightarrow \overrightarrow{GM}=(-10;11)$.\\
        Giả sử $\overrightarrow{GM} = h\overrightarrow{AB}+k\overrightarrow{AC} \Leftrightarrow \heva{&-10=2h-2k\\ &11=2h+6k} \Leftrightarrow \heva{&h= -\dfrac{19}{4}\\ &k= \dfrac{21}{4}.}$\\
        Vậy $\overrightarrow{GM} = -\dfrac{19}{4}\overrightarrow{AB}+\dfrac{21}{4}\overrightarrow{AC}$.
    \end{enumerate}
	}
\end{bt}

\begin{bt}%[0H9N1-5]
% nguon: AI soạn tương tự câu mẫu 42212c6c
% mau: 42212c6c
	Véc-tơ $\overrightarrow{a}=(0;-7)$ được phân tích theo hai véc-tơ đơn vị $\left(\overrightarrow{i};\overrightarrow{j}\right)$ như thế nào?
	\loigiai{
		Ta có $\overrightarrow{a}=(0;-7)\Rightarrow\overrightarrow{a}=0\overrightarrow{i}-7\overrightarrow{j}=-7\overrightarrow{j}$.
	}
\end{bt}

\begin{bt}%[0H9H1-5]
% nguon: AI soạn tương tự câu mẫu 42212c6c
% mau: 42212c6c nang_cao
	Trong mặt phẳng tọa độ $Oxy$, cho véc-tơ $\overrightarrow{b}=(3;-5)$. Biết $\overrightarrow{b}=m\overrightarrow{i}+n\overrightarrow{j}$ với $\overrightarrow{i},\overrightarrow{j}$ là hai véc-tơ đơn vị. Tính giá trị của biểu thức $S=m+n$.
	\loigiai{
		Ta có $\overrightarrow{b}=(3;-5)\Rightarrow\overrightarrow{b}=3\overrightarrow{i}-5\overrightarrow{j}$, suy ra $m=3$, $n=-5$.\\
		Vậy $S=m+n=3+(-5)=-2$.
	}
\end{bt}

\begin{bt}%[0H9H1-6]
% nguon: AI soạn tương tự câu mẫu b279dbb9
% mau: b279dbb9
Cho $\Delta ABC$ có $A(3; 7)$, $B(2;-2)$, $C(-6; 4)$. Tọa độ điểm $M(a;b)$ thuộc đoạn $BC$ sao cho $S_{\Delta MAB}=3 S_{\Delta MAC}$. Tính giá trị của biểu thức $a+b$.
	\loigiai{
		Kẻ $AH \perp BC$. $S_{\Delta MAB}=3 S_{\Delta MAC}\Leftrightarrow \dfrac{1}{2}AH\cdot MB=3\cdot \dfrac{1}{2}AH\cdot MC\Leftrightarrow MB=3MC$. \\
		Mà $\overrightarrow{MB}$ và $\overrightarrow{MC}$ ngược hướng $\Rightarrow \overrightarrow{MB}=-3\overrightarrow{MC}$ \\
		$\Leftrightarrow \heva{&x_B-x_M=-3(x_C-x_M)\\&y_B-y_M=-3(y_C-y_M)}\Leftrightarrow\heva{&2-x_M=-3(-6-x_M)\\&-2-y_M=-3(4-y_M)}$ \\
		$\Leftrightarrow\heva{&2-x_M=18+3x_M\\&-2-y_M=-12+3y_M}\Leftrightarrow\heva{&4x_M=-16\\&4y_M=10}\Leftrightarrow\heva{&x_M=-4\\&y_M=\dfrac{5}{2}}$. \\
		Suy ra $M\left(-4;\dfrac{5}{2}\right)$ nên $a+b=-4+\dfrac{5}{2}=-\dfrac{3}{2}$.
	}
\end{bt}

\begin{bt}%[0H9V1-6]
% nguon: AI soan tuong tu cau mau b279dbb9
% mau: b279dbb9 nang_cao
Cho $\Delta ABC$ có $A(6; 8)$, $B(5;-2)$, $C(-7; 4)$. Tọa độ điểm $M(a;b)$ thuộc đoạn $BC$ sao cho $S_{\Delta MAB}=4S_{\Delta MAC}$. Tính giá trị của biểu thức $a+b$.
\loigiai{
Kẻ $AH \perp BC$. $S_{\Delta MAB}=4S_{\Delta MAC}\Leftrightarrow \dfrac{1}{2}AH\cdot MB=4\cdot \dfrac{1}{2}AH\cdot MC\Leftrightarrow MB=4MC$.\\
Mà $\overrightarrow{MB}$ và $\overrightarrow{MC}$ ngược hướng $\Rightarrow \overrightarrow{MB}=-4\overrightarrow{MC}$\\
$\Leftrightarrow \heva{&x_B-x_M=-4(x_C-x_M)\\&y_B-y_M=-4(y_C-y_M)}\Leftrightarrow\heva{&5-x_M=-4(-7-x_M)\\&-2-y_M=-4(4-y_M)}$\\
Giải hệ: $5-x_M=28+4x_M\Leftrightarrow -5x_M=23\Leftrightarrow x_M=-\dfrac{23}{5}$;\\
$-2-y_M=-16+4y_M\Leftrightarrow -5y_M=-14\Leftrightarrow y_M=\dfrac{14}{5}$.\\
Suy ra $M\left(-\dfrac{23}{5};\dfrac{14}{5}\right)$.\\
Vậy $a+b=-\dfrac{23}{5}+\dfrac{14}{5}=-\dfrac{9}{5}$.
}
\end{bt}

\begin{bt}%[0H9N2-1]
% nguon: AI soạn tương tự câu mẫu 2dc15a56
% mau: 2dc15a56
	Cho hai véc-tơ $\overrightarrow{a}=(2;1),\overrightarrow{b}=(1;4)$.
	\begin{enumerate}
		\item Tìm toạ độ của véc-tơ $3\overrightarrow{a}+2\overrightarrow{b}$.
		\item Tính các tích vô hướng: $\overrightarrow{a} \cdot \overrightarrow{b},(2\overrightarrow{a}) \cdot (3\overrightarrow{b})$.
	\end{enumerate}
	\loigiai{
		\begin{enumerate}
			\item Ta có: $3\overrightarrow{a}+2\overrightarrow{b}=(3 \cdot 2+2 \cdot 1;3 \cdot 1+2 \cdot 4)=(8;11)$.\\
			\item Ta có: $\overrightarrow{a} \cdot \overrightarrow{b}=2 \cdot 1+1 \cdot 4=6$\\
			$2\overrightarrow{a}=(4;2)$ và $3\overrightarrow{b}=(3;12)$ nên $(2\overrightarrow{a}) \cdot (3\overrightarrow{b})=4 \cdot 3+2 \cdot 12=36$.
		\end{enumerate}
	}
\end{bt}

\begin{bt}%[0H9H2-1]
% nguon: AI soạn tương tự câu mẫu 2dc15a56
% mau: 2dc15a56 nang_cao
	Cho hai véc-tơ $\overrightarrow{a}=(2;-1),\overrightarrow{b}=(-1;3)$.
	\begin{enumerate}
		\item Tìm toạ độ của véc-tơ $3\overrightarrow{a}-2\overrightarrow{b}$.
		\item Tính các tích vô hướng: $\overrightarrow{a} \cdot \overrightarrow{b},(2\overrightarrow{a}) \cdot (3\overrightarrow{b})$.
	\end{enumerate}
	\loigiai{
		\begin{enumerate}
			\item Ta có: $3\overrightarrow{a}-2\overrightarrow{b}=(3 \cdot 2-2 \cdot (-1);3 \cdot (-1)-2 \cdot 3)=(8;-9)$.\\
			\item Ta có: $\overrightarrow{a} \cdot \overrightarrow{b}=2 \cdot (-1)+(-1) \cdot 3=-5$\\
			$2\overrightarrow{a}=(4;-2)$ và $3\overrightarrow{b}=(-3;9)$ nên $(2\overrightarrow{a}) \cdot (3\overrightarrow{b})=4 \cdot (-3)+(-2) \cdot 9=-30$.
		\end{enumerate}
	}
\end{bt}

\begin{bt}%[0H9N2-2]
% nguon: AI soạn tương tự câu mẫu 34f6d564
% mau: 34f6d564
	Cho $\overrightarrow{a}=(3; 2)$ và $\overrightarrow{b}=(-2;5)$. Tính $\cos \alpha$, với $\alpha =\left(\overrightarrow{a}, \overrightarrow{b}\right)$.
	\loigiai{
		Ta có $\cos \alpha =\dfrac{\overrightarrow{a}\cdot \overrightarrow{b}}{\left|\overrightarrow{a}\right|\cdot \left|\overrightarrow{b}\right|} =\dfrac{3\cdot (-2)+2\cdot 5}{\sqrt{3^2+2^2}\cdot \sqrt{(-2)^2+5^2}}=\dfrac{4}{\sqrt{377}}$.
	}
\end{bt}

\begin{bt}%[0H9H2-2]
% nguon: AI soạn tương tự câu mẫu 34f6d564
% mau: 34f6d564 nang_cao
	Cho $\overrightarrow{a}=(3; -1)$ và $\overrightarrow{b}=(2;4)$. Biết $\alpha =\left(\overrightarrow{a}, \overrightarrow{b}\right)$, tính $\cos \alpha$ và cho biết $\alpha$ là góc nhọn hay góc tù.
	\loigiai{
		Ta có $\cos \alpha =\dfrac{\overrightarrow{a}\cdot \overrightarrow{b}}{\left|\overrightarrow{a}\right|\cdot \left|\overrightarrow{b}\right|} =\dfrac{3\cdot 2+(-1)\cdot 4}{\sqrt{3^2+(-1)^2}\cdot \sqrt{2^2+4^2}}=\dfrac{2}{\sqrt{10}\cdot \sqrt{20}}=\dfrac{2}{\sqrt{200}}=\dfrac{2}{10\sqrt{2}}=\dfrac{\sqrt{2}}{10}$.\\
		Vì $\cos \alpha =\dfrac{\sqrt{2}}{10}>0$ nên $\alpha$ là góc nhọn.
	}
\end{bt}

\begin{bt}%[0H9H2-3]
% nguon: AI soạn tương tự câu mẫu a3548811
% mau: a3548811
	Trong mặt phẳng tọa độ $Oxy$, cho tam giác $ABC$ có $A(2;-1)$, $B(4;3)$ và $C(10;1)$. Tính chu vi tam giác $ABC$.
	\loigiai{
		Ta có $\heva{&\overrightarrow{AB}=(2;4)\\&\overrightarrow{AC}=(8;2)\\&\overrightarrow{BC}=(6;-2)} \Rightarrow \heva{&AB=\sqrt{2^2+4^2}=2\sqrt{5}\\&AC=\sqrt{8^2+2^2}=2\sqrt{17}\\&BC=\sqrt{6^2+(-2)^2}=2\sqrt{10}.}$\\
		Chu vi tam giác $ABC$ là $AB+AC+BC=2\sqrt{5}+2\sqrt{17}+2\sqrt{10}=2\sqrt{5}+2\sqrt{10}+2\sqrt{17}$.
	}
\end{bt}

\begin{bt}%[0H9V2-3]
% nguon: AI soạn tương tự câu mẫu a3548811
% mau: a3548811 nang_cao
	Trong mặt phẳng tọa độ $Oxy$, cho tam giác $ABC$ có $A(-2;3)$, $B(1;5)$ và $C(4;-1)$. Gọi $M$ là trung điểm của cạnh $BC$. Tính chu vi tam giác $ABM$.
	\loigiai{
		Vì $M$ là trung điểm của $BC$ nên $M\left(\dfrac{1+4}{2};\dfrac{5+(-1)}{2}\right)$, suy ra $M\left(\dfrac{5}{2};2\right)$.\\
		Ta có $\heva{&\overrightarrow{AB}=(3;2)\\&\overrightarrow{AM}=\left(\dfrac{9}{2};-1\right)\\&\overrightarrow{BM}=\left(\dfrac{3}{2};-3\right)} \Rightarrow \heva{&AB=\sqrt{3^2+2^2}=\sqrt{13}\\&AM=\sqrt{\left(\dfrac{9}{2}\right)^2+(-1)^2}=\dfrac{\sqrt{85}}{2}\\&BM=\sqrt{\left(\dfrac{3}{2}\right)^2+(-3)^2}=\dfrac{3\sqrt{5}}{2}.}$\\
		Chu vi tam giác $ABM$ là $AB+AM+BM=\sqrt{13}+\dfrac{\sqrt{85}}{2}+\dfrac{3\sqrt{5}}{2}=\sqrt{13}+\dfrac{\sqrt{85}+3\sqrt{5}}{2}$.
	}
\end{bt}

\begin{bt}%[0H9N2-4]
% nguon: AI soạn tương tự câu mẫu 0e8eb5bf
% mau: 0e8eb5bf
Trong mặt phẳng tọa độ $Oxy$, cho điểm $B(2;-3)$, đường thẳng $\Delta_1\colon 2x-y+5=0$ và đường thẳng $\Delta_2\colon x+3y-20=0$.
	\begin{enumerate}
		\item Tính góc giữa đường thẳng $\Delta_1$ và $\Delta_2$.
		\item Tính khoảng cách từ điểm $B$ đến đường thẳng $\Delta_2$.
	\end{enumerate}
	\loigiai{
		\begin{enumerate}
			\item $\Delta_1$ có véc-tơ pháp tuyến $\overrightarrow{n}_1=(2;-1)$ và $\Delta_2$ có véc-tơ pháp tuyến $\overrightarrow{n}_2=(1;3)$.\\
			$\cos(\Delta_1;\Delta_2)=\left|\cos(\overrightarrow{n}_1;\overrightarrow{n}_2)\right|=\dfrac{|2\cdot1+(-1)\cdot3|}{\sqrt{2^2+(-1)^2}\cdot\sqrt{1^2+3^2}}=\dfrac{1}{\sqrt{5}\cdot\sqrt{10}}=\dfrac{\sqrt{2}}{10}$.\\
			Suy ra $(\Delta_1;\Delta_2)\approx 81{,}9^\circ$.
			\item $\mathrm{d}(B,\Delta_2)=\dfrac{|1\cdot2+3\cdot(-3)-20|}{\sqrt{1^2+3^2}}=\dfrac{27}{\sqrt{10}}=\dfrac{27\sqrt{10}}{10}$.
		\end{enumerate}
	}
\end{bt}

\begin{bt}%[0H9H2-4]
% nguon: AI soạn tương tự câu mẫu 0e8eb5bf
% mau: 0e8eb5bf nang_cao
Trong mặt phẳng tọa độ $Oxy$, cho điểm $M(4;-2)$, đường thẳng $\Delta_1\colon 3x-y+10=0$ và đường thẳng $\Delta_2\colon x+2y-30=0$.
	\begin{enumerate}
		\item Tính góc giữa đường thẳng $\Delta_1$ và $\Delta_2$.
		\item Tính khoảng cách từ điểm $M$ đến đường thẳng $\Delta_2$, biết rằng điểm $M$ là trung điểm của đoạn thẳng $NP$ với $N(1;3)$ và $P$ có tung độ bằng $-7$. Xác định tọa độ điểm $P$ rồi tính khoảng cách từ $P$ đến $\Delta_2$.
	\end{enumerate}
	\loigiai{
		\begin{enumerate}
			\item $\Delta_1$ có véc-tơ pháp tuyến $\overrightarrow{n}_1=(3;-1)$ và $\Delta_2$ có véc-tơ pháp tuyến $\overrightarrow{n}_2=(1;2)$.\\
			$\cos(\Delta_1;\Delta_2)=\left|\cos(\overrightarrow{n}_1;\overrightarrow{n}_2)\right|=\dfrac{|3\cdot1+(-1)\cdot2|}{\sqrt{3^2+(-1)^2}\cdot\sqrt{1^2+2^2}}=\dfrac{1}{\sqrt{10}\cdot\sqrt{5}}=\dfrac{\sqrt{2}}{10}$.\\
			Suy ra $(\Delta_1;\Delta_2)\approx 81{,}87^\circ$.
			\item Gọi $P(x_P;-7)$. Vì $M(4;-2)$ là trung điểm của $NP$ với $N(1;3)$ nên
			$x_P=2x_M-x_N=2\cdot4-1=7$ và $y_P=2y_M-y_N=2\cdot(-2)-3=-7$ (thỏa mãn).\\
			Vậy $P(7;-7)$.\\
			$\mathrm{d}(P,\Delta_2)=\dfrac{|7+2\cdot(-7)-30|}{\sqrt{1^2+2^2}}=\dfrac{|7-14-30|}{\sqrt{5}}=\dfrac{37}{\sqrt{5}}=\dfrac{37\sqrt{5}}{5}$.
		\end{enumerate}
	}
\end{bt}
